% فصل چهارم: پیاده‌سازی و ارزیابی نتایج
\chapter{پیاده‌سازی و ارزیابی نتایج}
\label{ch:implementation}

در این فصل، جزئیات محاسباتی و نتایج تجربی حاصل از پیاده‌سازی چارچوب شبکه‌های عصبی آگاه از فیزیک ارائه می‌گردد. ابتدا بستر نرم‌افزاری و ابزارهای مورد استفاده معرفی شده و سپس حل‌گر مرجع طیفی توسعه‌یافته برای معادله غیرخطی برگرز تشریح و با آزمون‌های مستقل اعتبارسنجی می‌شود. در ادامه، نتایج حل مستقیم معادله با مدل زمان‌پیوسته و نتایج کشف پارامترها در مسئله معکوس در سناریوهای داده تمیز و نویزی گزارش و تحلیل می‌گردند. در پایان، یافته‌های تجربی با نتایج مقاله مرجع مقایسه و نقاط قوت و محدودیت‌های روش تبیین می‌شود.

\section{بستر محاسباتی و ابزارهای پیاده‌سازی}
\label{sec:env}

کلیه الگوریتم‌های محاسباتی این پژوهش با زبان برنامه‌نویسی پایتون نسخه ۳.۱۱ پیاده‌سازی شده‌اند. برای پیاده‌سازی ساختار شبکه عصبی، کامپایل توابع و انجام مشتق‌گیری خودکار، کتابخانه جکس\LTRfootnote{JAX} به کار گرفته شده است \cite{jax2018}. جکس امکان ترکیب تبدیل‌های تابعی نظیر مشتق‌گیری خودکار (\lr{grad})، برداری‌سازی خودکار (\lr{vmap}) و کامپایل در زمان اجرا (\lr{jit}) را بر بستر کامپایلر \lr{XLA} فراهم می‌سازد. انجام محاسبات عددی و ماتریسی با استفاده از کتابخانه‌های نامپای و سای‌پای، و رسم نمودارها با کتابخانه مت‌پلات‌لیب صورت پذیرفته است. الگوریتم بهینه‌سازی ترکیبی شامل پیاده‌سازی مستقیم بهینه‌ساز آدام و استفاده از رابط الگوریتم بهینه‌سازی مقید \lr{L-BFGS-B} در کتابخانه سای‌پای است.

تمامی مراحل آموزش و آزمایش‌های عددی بر روی واحد پردازش مرکزی (\lr{CPU}) و با دقت ممیز شناور ۶۴ بیتی (\lr{Float64}) اجرا شده‌اند تا امکان بازتولید دقیق نتایج بر روی تجهیزات استاندارد دانشگاهی فراهم باشد. به‌منظور تضمین تکرارپذیری، بذرهای تولید اعداد تصادفی ثابت نگاه داشته شده و کد کامل برنامه به همراه راهنمای اجرا در پیوست الف و پوشه اختصاصی \lr{code} ارائه گردیده است.

\section{حل‌گر مرجع طیفی و اعتبارسنجی آن}
\label{sec:reference}

جهت ارزیابی مستقل و دقیق عملکرد شبکه آگاه از فیزیک، یک حل‌گر عددی مرجع مستقل برای معادله برگرز با لزجت $\nu = 0.01/\pi$ توسعه داده شد. در این حل‌گر، گسسته‌سازی مکانی با روش طیفی فوریه بر روی ۱۰۲۴ مود با اعمال قاعده پالایش دو-سوم برای حذف خطای هم‌پوشانی\LTRfootnote{Aliasing Error} صورت گرفته است \cite{trefethen2000}. گام‌زنی زمانی با استفاده از روش تفاضل زمانی نمایی رانگه-کوتای مرتبه چهار (\lr{ETDRK4}) با طول گام زمانی $\Delta t = 0.001$ انجام می‌پذیرد که برای معادلات دارای جملات پخش سفت، پایداری فوق‌العاده‌ای ارائه می‌دهد \cite{kassam2005}. ماتریس پاسخ حاصل بر روی یک شبکه زمانی-مکانی با ابعاد $101 \times 1024$ ذخیره شده و به‌عنوان معیار حل دقیق در ارزیابی‌ها استفاده شده است.

صحت عملکرد حل‌گر مرجع با دو آزمون مستقل اثبات گردید:
\begin{enumerate}
\item \textbf{آزمون همگرایی عددی:} با نصف کردن گام زمانی ($\Delta t = 0.0005$)، بیشینه تغییر نسبی جواب در مرتبه $6.2\times 10^{-9}$ و با دو برابر کردن تفکیک مکانی ($N_x=2048$)، تغییر نسبی برابر $4.1\times 10^{-4}$ ثبت گردید که مؤید همگرایی مطلوب روش است.
\item \textbf{آزمون مقایسه با جواب نیمه‌تحلیلی کول-هوپف:} جواب معادله با تبدیل تحلیلی کول-هوپف \cite{hopf1950} از طریق انتگرال‌گیری عددی مستقیم با دقت بسیار بالا محاسبه شد. بیشترین اختلاف در ۱۹ نقطه نمونه‌برداری در زمان نهایی $t=1$ کمتر از $6\times 10^{-4}$ بوده و در نواحی دور از جبهه شوک به مرتبه $10^{-8}$ کاهش یافت. همچنین مطابقت با بسط‌های تحلیلی سری فوریه \cite{benton1972} نیز بررسی گردید.
\end{enumerate}

این آزمون‌ها نشان می‌دهند که دقت حل مرجع چندین مرتبه بزرگی فراتر از دقت هدف شبکه‌های عصبی بوده و مبنایی کاملاً قابل‌اعتماد برای ارزیابی خطاست.

\section{حل مستقیم معادله برگرز}
\label{sec:forward}

در مسئله مستقیم، هدف محاسبه جواب $u(t,x)$ در دامنه مکانی $x \in [-1,1]$ و زمانی $t \in [0,1]$ با شرط اولیه $u(0,x) = -\sin(\pi x)$ و شرایط مرزی همگن دیریکله $u(t,-1) = u(t,1) = 0$ و ضرایب معلوم $\lambda_1 = 1$ و $\lambda_2 = 0.01/\pi$ است. ساختار شبکه عصبی مورد استفاده، پرسپترون چندلایه با ۴ لایه پنهان، ۵۰ نورون در هر لایه و تابع فعال‌ساز $\tanh$ است که در مجموع شامل ۷۸۵۱ پارامتر قابل آموزش می‌باشد.

داده‌های آموزشی روی مرزها و شرط اولیه شامل ۲۰۰ نقطه ($N_u = 200$) شامل ۱۰۰ نقطه در $t=0$ و ۱۰۰ نقطه روی مرزهای مکانی است. تعداد نقاط هم‌مکانی درون دامنه برابر ۱۰۰۰۰ نقطه ($N_f = 10000$) انتخاب شد که نیمی از آن‌ها به‌صورت تصادفی یکنواخت و نیمی دیگر با توزیع نرمال متمرکز حول خط $x=0$ (ناحیه تشکیل جبهه شوک) نمونه‌برداری گردیدند. موقعیت نقاط آموزشی در شکل \ref{fig:fwdpoints} و مشخصات ساختاری آموزش در جدول \ref{tab:fwdconfig} خلاصه شده است.

\begin{figure}[t]
\centering
\includegraphics[width=0.85\textwidth]{forward_points}
\caption{توزیع نقاط آموزشی در مسئله مستقیم معادله برگرز: نقاط داده اولیه و مرزی ($N_u=200$) و نقاط هم‌مکانی باقیمانده ($N_f=10000$)}
\label{fig:fwdpoints}
\end{figure}

\begin{table}[t]
\caption{تنظیمات ساختاری و پارامترهای آموزش مدل زمان‌پیوسته در مسئله مستقیم}
\label{tab:fwdconfig}
\centering
\singlespacing
\begin{tabular}{|p{4.8cm}|p{8.5cm}|}
\hline
\textbf{مولفه ساختاری} & \textbf{مشخصات انتخابی} \\
\hline
معماری شبکه عصبی & ۴ لایه پنهان، ۵۰ نورون در هر لایه، فعال‌ساز تانژانت هیپربولیک \\
\hline
تعداد کل پارامترهای شبکه & ۷۸۵۱ پارامتر وزن و بایاس \\
\hline
نقاط داده اولیه و مرزی ($N_u$) & ۲۰۰ نقطه (۱۰۰ شرط اولیه و ۱۰۰ شرط مرزی) \\
\hline
نقاط هم‌مکانی باقیمانده ($N_f$) & ۱۰۰۰۰ نقطه (۵۰۰۰ یکنواخت و ۵۰۰۰ متمرکز در ناحیه شوک) \\
\hline
الگوریتم مقداردهی اولیه & روش گزاویه (گلوروت) \\
\hline
راهبرد بهینه‌سازی & ۶۰۰۰ گام آدام ($\eta=10^{-3}$) و ادامه با \lr{L-BFGS-B} تا ۶۰۰۰ فراخوانی \\
\hline
دقت محاسباتی ممیز شناور & ۶۴ بیتی (\lr{Float64}) \\
\hline
\end{tabular}
\end{table}

فرایند آموزش طی دو مرحله متوالی اجرا گردید: ۶۰۰۰ گام با الگوریتم آدام با نرخ یادگیری $10^{-3}$ و سپس ادامه بهینه‌سازی با روش \lr{L-BFGS-B} تا دستیابی به همگرایی. مقدار تابع هزینه نهایی در انتهای مرحله آدام برابر \FwdLossAdam\ و در پایان مرحله \lr{L-BFGS} برابر \FwdLossLBFGS\ حاصل شد که در آن سهم مولفه خطای داده $\mathrm{MSE}_u$ برابر \FwdMSEu\ و مولفه خطای باقیمانده فیزیک $\mathrm{MSE}_f$ برابر \FwdMSEf\ بود. تاریخچه تغییرات تابع هزینه در شکل \ref{fig:loss} رسم گردیده است.

\begin{figure}[t]
\centering
\includegraphics[width=0.85\textwidth]{loss_history}
\caption{تاریخچه کاهش تابع هزینه در مسئله مستقیم: همگرایی مرحله اول با الگوریتم آدام و اصلاح نهایی با روش شبه‌نیوتنی \lr{L-BFGS}}
\label{fig:loss}
\end{figure}

ارزیابی شبکه آموزش‌دیده بر روی کل شبکه مرجع با ۱۰۳۴۲۴ نقطه ($101 \times 1024$) نشان‌دهنده دستیابی به خطای نسبی در نرم $L_2$ برابر \FwdRelErr\ است:
\begin{equation}
\text{Relative } L_2 \text{ Error} = \frac{\|u_{\text{pred}} - u_{\text{ref}}\|_2}{\|u_{\text{ref}}\|_2} = 2.15\times 10^{-2}.
\end{equation}

مقایسه پروفیل‌های زمانی پاسخ پیش‌بینی‌شده با حل مرجع در مقاطع $t=0.25$، $t=0.5$ و $t=0.75$ در شکل \ref{fig:snapshots} نمایش داده شده است. شکل \ref{fig:errmap} نقشه خطای قدرمطلق را در سراسر دامنه زمانی-مکانی نشان می‌دهد. این نمودار تبیین می‌کند که خطای تقریب در بخش اعظم دامنه ناچیز بوده و بیشینه آن در لایه باریک گرادیان تند شوک (حوالی $x=0$ برای $t > 0.3$) متمرکز است.

\begin{figure}[t]
\centering
\includegraphics[width=0.85\textwidth]{burgers_snapshots}
\caption{مقایسه پروفیل‌های زمانی حل دقیق مرجع و پیش‌بینی شبکه آگاه از فیزیک در سه لحظه $t=0.25, 0.5, 0.75$}
\label{fig:snapshots}
\end{figure}

\begin{figure}[t]
\centering
\includegraphics[width=0.85\textwidth]{error_heatmap}
\caption{نقشه توزیع خطای قدرمطلق $|u_{\text{ref}} - u_{\text{PINN}}|$ در دامنه زمانی-مکانی در مسئله مستقیم}
\label{fig:errmap}
\end{figure}

جدول \ref{tab:comparison} مقایسه نتایج این پژوهش را با مقاله مرجع \cite{raissi2019} نشان می‌دهد. تفاوت اندک در دقت به عواملی نظیر تفاوت در تعداد لایه‌ها و عمق شبکه (۴ لایه ۵۰ نورونی در برابر ۹ لایه ۲۰ نورونی)، تفاوت پلتفرم سخت‌افزاری پردازش و سقف تعداد تکرارهای بهینه‌سازی مربوط می‌شود.

\begin{table}[t]
\caption{مقایسه تطبیقی نتایج حل مستقیم با مقاله مرجع رئیسی و همکاران (۲۰۱۹)}
\label{tab:comparison}
\centering
\singlespacing
\begin{tabular}{|p{3.2cm}|p{3.2cm}|p{2.3cm}|p{2.3cm}|p{2.4cm}|}
\hline
\textbf{منبع پژوهش} & \textbf{معماری شبکه} & \textbf{داده مرزی} & \textbf{نقاط هم‌مکانی} & \textbf{خطای نسبی $L_2$} \\
\hline
رئیسی و همکاران \cite{raissi2019} & ۹ لایه، ۲۰ نورون & ۱۰۰ & ۱۰۰۰۰ & $6.7\times 10^{-4}$ \\
\hline
پژوهش حاضر & ۴ لایه، ۵۰ نورون & ۲۰۰ & ۱۰۰۰۰ & \FwdRelErr \\
\hline
\end{tabular}
\end{table}

\section{کشف پارامترهای نامعلوم در مسئله معکوس}
\label{sec:inverse-exp}

در مسئله معکوس، فرض بر این است که ضرایب معادله برگرز در عملگر باقیمانده $f = u_t + \lambda_1 u u_x - \lambda_2 u_{xx}$ مجهول بوده و هدف، استخراج هم‌زمان پارامترهای $\lambda_1$ و $\lambda_2$ از روی مشاهدات تجربی پراکنده است. مجموعه داده‌های آموزشی شامل ۵۰۰۰ مشاهده پراکنده ($N=5000$) در دامنه زمانی-مکانی است که نیمی به‌صورت یکنواخت و نیمی متمرکز حول $x=0$ توزیع شده‌اند (شکل \ref{fig:invpoints}).

ساختار شبکه شامل ۵ لایه پنهان با ۳۰ نورون در هر لایه است. مقادیر اولیه هر دو ضریب برابر صفر ($\lambda_1=0, \lambda_2=0$) تعیین شد. بهینه‌سازی مطابق راهبرد ترکیبی ۶۰۰۰ گام آدام و ادامه با \lr{L-BFGS-B} انجام گرفت.

\begin{figure}[t]
\centering
\includegraphics[width=0.85\textwidth]{inverse_points}
\caption{توزیع مکانی-زمانی ۵۰۰۰ نقطه مشاهده پراکنده در مسئله معکوس؛ طیف رنگی نمایانگر مقدار متغیر حالت $u(t,x)$ است}
\label{fig:invpoints}
\end{figure}

مسئله معکوس در دو سناریوی آزمایشی بررسی شد:
\begin{enumerate}
\item \textbf{سناریوی داده تمیز (بدون نویز):} داده‌ها مستقیماً از حل مرجع استخراج شدند.
\item \textbf{سناریوی داده نویزی:} داده‌ها با نویز گاوسی با انحراف معیار برابر یک درصد انحراف معیار سیگنال آلوده شدند ($u_{\text{noisy}} = u + 0.01 \times \mathrm{std}(u) \times \mathcal{N}(0,1)$).
\end{enumerate}

جدول \ref{tab:inverse} مقادیر ضرایب شناسایی‌شده و درصد خطای نسبی آن‌ها را در مقایسه با مقادیر واقعی ($\lambda_1^* = 1.0$ و $\lambda_2^* = 0.01/\pi \approx 0.003183$) گزارش می‌نماید.

\begin{table}[t]
\caption{نتایج تخمین پارامترهای مجهول معادله برگرز در مسئله معکوس}
\label{tab:inverse}
\centering
\singlespacing
\begin{tabular}{|p{2.8cm}|p{2.6cm}|p{2.4cm}|p{2.6cm}|p{2.4cm}|}
\hline
\textbf{سناریوی داده} & \textbf{$\lambda_1$ برآوردی} & \textbf{خطای $\lambda_1$} & \textbf{$\lambda_2$ برآوردی} & \textbf{خطای $\lambda_2$} \\
\hline
بدون نویز (تمیز) & \LamIClean & \LamIErrClean & \LamIIClean & \LamIIErrClean \\
\hline
آلوده به ۱٪ نویز & \LamINoisy & \LamIErrNoisy & \LamIINoisy & \LamIIErrNoisy \\
\hline
مقادیر واقعی هدف & \multicolumn{2}{|c|}{$1.000$} & \multicolumn{2}{|c|}{$0.01/\pi \approx 0.003183$} \\
\hline
\end{tabular}
\end{table}

\section{بحث و تحلیل نتایج}
\label{sec:discussion}

تحلیل نتایج عددی حاصل از پیاده‌سازی، نکات علمی و مهندسی متعددی را در خصوص رفتار شبکه‌های آگاه از فیزیک روشن می‌سازد:

\textbf{کارایی داده‌ای بالا.} در مسئله مستقیم، تقریب میدان پیوسته پاسخ در بیش از یکصد هزار نقطه محاسباتی تنها بر مبنای ۲۰۰ نقطه داده اولیه و مرزی صورت پذیرفت. این نسبت نشان‌دهنده نقش برجسته قید باقیمانده فیزیکی در کاهش ابعاد فضای جست‌وجو و جلوگیری از بیش‌برازش است.

\textbf{چالش لایه‌های با گرادیان تند.} تمرکز عمده خطا در حوالی جبهه شوک نشان می‌دهد که پدیده سلبیت و تغییرات ناگهانی گرادیان، اصلی‌ترین گلوگاه دقتی شبکه‌های عصبی استاندارد در حل معادلات هذلولوی و شبه‌هذلولوی است. نمونه‌برداری غیریکنواخت و متمرکز در این ناحیه نقش تعیین‌کننده‌ای در کنترل خطای عددی ایفا نمود.

\textbf{تحلیل حساسیت و خطای شناسایی ضریب پخش.} خطای شناسایی ضریب انتقال ($\lambda_1$) در هر دو حالت کمتر از ۹ درصد بوده است، در حالی که خطای ضریب پخش ($\lambda_2$) در بازه ۲۵ تا ۳۳ درصد قرار دارد. این تفاوت ناشی از ماهیت فیزیکی مسئله است: با توجه به کوچکی مقدار $\lambda_2 \approx 0.00318$، اثر جمله پخشی $\lambda_2 u_{xx}$ در غالب نقاط دامنه در برابر جمله جابه‌جایی ناچیز است و صرفاً در لایه بسیار باریک شوک اثرگذار می‌شود؛ بنابراین گرادیان اطلاعاتی موجود در داده‌ها برای تنظیم $\lambda_2$ بسیار ضعیف‌تر از $\lambda_1$ است.

\textbf{پایداری در برابر نویز و رفتار بهینه‌سازی.} نتایج جدول \ref{tab:inverse} نشان می‌دهد که افزودن نویز یک درصد ساختار تخمین پارامترها را دچار واگرایی نکرده و خطای شناسایی در همان مرتبه داده‌های تمیز باقی مانده است. تفاوت جزئی در خطای عددی میان دو سناریو، ناشی از نوسانات تصادفی بهینه‌سازی غیرمحدب و تفاوت در بذر تصادفی مقداردهی اولیه است و بیانگر پایداری رضایت‌بخش روش در برابر نویز به دلیل فیلترسازی طبیعی توسط شبکه عصبی هموار است.

\textbf{مقایسه با روش‌های کلاسیک.} در حالی که حل‌گرهای کلاسیک طیفی در تولید پاسخ برای یک مسئله معین با سرعت و دقت بسیار بالاتری عمل می‌کنند، چارچوب شبکه‌های آگاه از فیزیک در مسائلی که داده‌ها ناقص بوده، شرایط مرزی ناشناخته باشند، یا پارامترهای سیستم مجهول باشند، انعطاف‌پذیری بی‌نظیری از خود نشان می‌دهد.

\textbf{محدودیت‌های پژوهش حاضر.} این پژوهش بر روی معادله برگرز یک‌بعدی و مدل زمان‌پیوسته متمرکز بوده و بهینه‌سازی ابرپارامترها به‌دلیل محدودیت توان پردازشی در مقیاس محدودی انجام گرفته است. همچنین تحلیل عدم‌قطعیت پارامترهای استخراج‌شده بررسی نشده است.

\section{جمع‌بندی فصل}
\label{sec:ch4-summary}

در این فصل، پیاده‌سازی و ارزیابی محاسباتی مدل زمان‌پیوسته شبکه آگاه از فیزیک بر روی معادله برگرز ارائه شد. حل‌گر مرجع طیفی با گام‌زنی نمایی مرتبه چهار توسعه یافته و اعتبارسنجی شد. نتایج نشان داد که مدل زمان‌پیوسته با ۲۰۰ نقطه داده اولیه-مرزی به خطای نسبی \FwdRelErr\ در مسئله مستقیم دست می‌یابد. در مسئله معکوس نیز ضرایب با دقت مناسبی استخراج شدند و پایداری روش در حضور نویز داده‌ها مورد تایید قرار گرفت.
